% test-008.tex -- tagged-PDF test de \spnord
%
% Compilar:  lualatex-dev test-008.tex
% Validar:   show-pdf-tags --xml test-008.pdf | rnv-wrapp
%            verapdf --flavour ua2 --format text test-008.pdf
%
% Test de prueba de \spnord: ordinales en modo matematico, con sus
% llaves small-caps y read-arg. Los "doc:" son las lecturas que se esperan
% de MathCAT con la regla de ordinales (#999), aun sin comprobar en NVDA.
\DocumentMetadata
  {
    lang=es-CL,
    pdfversion=2.0,
    pdfstandard=ua-2,
    tagging=on,
    tagging-setup={math/setup={mathml-SE}},
  }
\documentclass{article}
\usepackage[spanish]{babel}
\usepackage{unicode-math}
\usepackage{spanish-mathml-intent}
\usepackage{hyperref}
\hypersetup
  {
    pdfdisplaydoctitle = true,
    pdftitle           = {test-008: ordinales con spnord},
  }

\begin{document}

\section*{A. Ordinales en una fórmula}

% doc: «primero»
1. Masculino: $\spnord{1.o}$.

% doc: «primera»
2. Femenino: $\spnord{1.a}$.

% doc: «tercer»
3. Apocopado: $\spnord{3.er}$.

% doc: «quintas»
4. Femenino plural: $\spnord{5.as}$.

% doc: «vigésima primera»
5. Compuesto: $\spnord{21.a}$.

% doc: «segundo»
6. Con la letra º: $\spnord{2.º}$.

% doc: «segundo»
7. Con versalitas: $\spnord[small-caps=true]{2.o}$.

% doc: «octavo»
8. Con lectura propia: $\spnord[read-arg={octavo}]{8.^{\text{avo}}}$.


\end{document}
